\section{}

\paragraph{1234年}\eventanchor{SS-1234-REGNAL-YEAR}{SS-1234-REGNAL-YEAR}　南宋以宋理宗在位第10年作为诏令、赋役与外交文书的纪年框架。账本所列《宋史》卷24—47〈本纪〉；《建炎以来系年要录》本年条；《宋会要辑稿》相关门类为本条来源起点；该条可固定编年位置，具体执行、规模和地方社会影响仍须保留来源差异。

\paragraph{1235年}\eventanchor{SS-1235-EVENT-139}{SS-1235-EVENT-139}　蒙古军开始多方向进攻南宋边境。账本所列《宋史》卷24—47〈本纪〉；《建炎以来系年要录》本年条；《宋会要辑稿》相关门类为本条来源起点；该条可固定编年位置，具体执行、规模和地方社会影响仍须保留来源差异。

\paragraph{1235年}\eventanchor{SS-1235-REGNAL-YEAR}{SS-1235-REGNAL-YEAR}　南宋以宋理宗在位第11年作为诏令、赋役与外交文书的纪年框架。账本所列《宋史》卷24—47〈本纪〉；《建炎以来系年要录》本年条；《宋会要辑稿》相关门类为本条来源起点；该条可固定编年位置，具体执行、规模和地方社会影响仍须保留来源差异。

\paragraph{1236年}\eventanchor{SS-1236-REGNAL-YEAR}{SS-1236-REGNAL-YEAR}　南宋以宋理宗在位第12年作为诏令、赋役与外交文书的纪年框架。账本所列《宋史》卷24—47〈本纪〉；《建炎以来系年要录》本年条；《宋会要辑稿》相关门类为本条来源起点；该条可固定编年位置，具体执行、规模和地方社会影响仍须保留来源差异。

\paragraph{1237年}\eventanchor{SS-1237-REGNAL-YEAR}{SS-1237-REGNAL-YEAR}　南宋以宋理宗在位第13年作为诏令、赋役与外交文书的纪年框架。账本所列《宋史》卷24—47〈本纪〉；《建炎以来系年要录》本年条；《宋会要辑稿》相关门类为本条来源起点；该条可固定编年位置，具体执行、规模和地方社会影响仍须保留来源差异。

\paragraph{1238年}\eventanchor{SS-1238-REGNAL-YEAR}{SS-1238-REGNAL-YEAR}　南宋以宋理宗在位第14年作为诏令、赋役与外交文书的纪年框架。账本所列《宋史》卷24—47〈本纪〉；《建炎以来系年要录》本年条；《宋会要辑稿》相关门类为本条来源起点；该条可固定编年位置，具体执行、规模和地方社会影响仍须保留来源差异。

\paragraph{1239年}\eventanchor{SS-1239-REGNAL-YEAR}{SS-1239-REGNAL-YEAR}　南宋以宋理宗在位第15年作为诏令、赋役与外交文书的纪年框架。账本所列《宋史》卷24—47〈本纪〉；《建炎以来系年要录》本年条；《宋会要辑稿》相关门类为本条来源起点；该条可固定编年位置，具体执行、规模和地方社会影响仍须保留来源差异。

\paragraph{1240年}\eventanchor{SS-1240-REGNAL-YEAR}{SS-1240-REGNAL-YEAR}　南宋以宋理宗在位第16年作为诏令、赋役与外交文书的纪年框架。账本所列《宋史》卷24—47〈本纪〉；《建炎以来系年要录》本年条；《宋会要辑稿》相关门类为本条来源起点；该条可固定编年位置，具体执行、规模和地方社会影响仍须保留来源差异。

\paragraph{1241年}\eventanchor{SS-1241-REGNAL-YEAR}{SS-1241-REGNAL-YEAR}　南宋以宋理宗在位第17年作为诏令、赋役与外交文书的纪年框架。账本所列《宋史》卷24—47〈本纪〉；《建炎以来系年要录》本年条；《宋会要辑稿》相关门类为本条来源起点；该条可固定编年位置，具体执行、规模和地方社会影响仍须保留来源差异。

\paragraph{1242年}\eventanchor{SS-1242-REGNAL-YEAR}{SS-1242-REGNAL-YEAR}　南宋以宋理宗在位第18年作为诏令、赋役与外交文书的纪年框架。账本所列《宋史》卷24—47〈本纪〉；《建炎以来系年要录》本年条；《宋会要辑稿》相关门类为本条来源起点；该条可固定编年位置，具体执行、规模和地方社会影响仍须保留来源差异。

\paragraph{1243年}\eventanchor{SS-1243-REGNAL-YEAR}{SS-1243-REGNAL-YEAR}　南宋以宋理宗在位第19年作为诏令、赋役与外交文书的纪年框架。账本所列《宋史》卷24—47〈本纪〉；《建炎以来系年要录》本年条；《宋会要辑稿》相关门类为本条来源起点；该条可固定编年位置，具体执行、规模和地方社会影响仍须保留来源差异。

\paragraph{1244年}\eventanchor{SS-1244-REGNAL-YEAR}{SS-1244-REGNAL-YEAR}　南宋以宋理宗在位第20年作为诏令、赋役与外交文书的纪年框架。账本所列《宋史》卷24—47〈本纪〉；《建炎以来系年要录》本年条；《宋会要辑稿》相关门类为本条来源起点；该条可固定编年位置，具体执行、规模和地方社会影响仍须保留来源差异。

\paragraph{1245年}\eventanchor{SS-1245-REGNAL-YEAR}{SS-1245-REGNAL-YEAR}　南宋以宋理宗在位第21年作为诏令、赋役与外交文书的纪年框架。账本所列《宋史》卷24—47〈本纪〉；《建炎以来系年要录》本年条；《宋会要辑稿》相关门类为本条来源起点；该条可固定编年位置，具体执行、规模和地方社会影响仍须保留来源差异。

\paragraph{1246年}\eventanchor{SS-1246-REGNAL-YEAR}{SS-1246-REGNAL-YEAR}　南宋以宋理宗在位第22年作为诏令、赋役与外交文书的纪年框架。账本所列《宋史》卷24—47〈本纪〉；《建炎以来系年要录》本年条；《宋会要辑稿》相关门类为本条来源起点；该条可固定编年位置，具体执行、规模和地方社会影响仍须保留来源差异。

\paragraph{1247年}\eventanchor{SS-1247-REGNAL-YEAR}{SS-1247-REGNAL-YEAR}　南宋以宋理宗在位第23年作为诏令、赋役与外交文书的纪年框架。账本所列《宋史》卷24—47〈本纪〉；《建炎以来系年要录》本年条；《宋会要辑稿》相关门类为本条来源起点；该条可固定编年位置，具体执行、规模和地方社会影响仍须保留来源差异。

\paragraph{1248年}\eventanchor{SS-1248-REGNAL-YEAR}{SS-1248-REGNAL-YEAR}　南宋以宋理宗在位第24年作为诏令、赋役与外交文书的纪年框架。账本所列《宋史》卷24—47〈本纪〉；《建炎以来系年要录》本年条；《宋会要辑稿》相关门类为本条来源起点；该条可固定编年位置，具体执行、规模和地方社会影响仍须保留来源差异。

\paragraph{1249年}\eventanchor{SS-1249-REGNAL-YEAR}{SS-1249-REGNAL-YEAR}　南宋以宋理宗在位第25年作为诏令、赋役与外交文书的纪年框架。账本所列《宋史》卷24—47〈本纪〉；《建炎以来系年要录》本年条；《宋会要辑稿》相关门类为本条来源起点；该条可固定编年位置，具体执行、规模和地方社会影响仍须保留来源差异。

\paragraph{1250年}\eventanchor{SS-1250-REGNAL-YEAR}{SS-1250-REGNAL-YEAR}　南宋以宋理宗在位第26年作为诏令、赋役与外交文书的纪年框架。账本所列《宋史》卷24—47〈本纪〉；《建炎以来系年要录》本年条；《宋会要辑稿》相关门类为本条来源起点；该条可固定编年位置，具体执行、规模和地方社会影响仍须保留来源差异。

\paragraph{1251年}\eventanchor{SS-1251-REGNAL-YEAR}{SS-1251-REGNAL-YEAR}　南宋以宋理宗在位第27年作为诏令、赋役与外交文书的纪年框架。账本所列《宋史》卷24—47〈本纪〉；《建炎以来系年要录》本年条；《宋会要辑稿》相关门类为本条来源起点；该条可固定编年位置，具体执行、规模和地方社会影响仍须保留来源差异。

\paragraph{1252年}\eventanchor{SS-1252-REGNAL-YEAR}{SS-1252-REGNAL-YEAR}　南宋以宋理宗在位第28年作为诏令、赋役与外交文书的纪年框架。账本所列《宋史》卷24—47〈本纪〉；《建炎以来系年要录》本年条；《宋会要辑稿》相关门类为本条来源起点；该条可固定编年位置，具体执行、规模和地方社会影响仍须保留来源差异。

\paragraph{1253年}\eventanchor{SS-1253-EVENT-140}{SS-1253-EVENT-140}　蒙古征服大理，南宋西南战略环境改变。账本所列《宋史》卷24—47〈本纪〉；《建炎以来系年要录》本年条；《宋会要辑稿》相关门类为本条来源起点；该条可固定编年位置，具体执行、规模和地方社会影响仍须保留来源差异。

\paragraph{1253年}\eventanchor{SS-1253-REGNAL-YEAR}{SS-1253-REGNAL-YEAR}　南宋以宋理宗在位第29年作为诏令、赋役与外交文书的纪年框架。账本所列《宋史》卷24—47〈本纪〉；《建炎以来系年要录》本年条；《宋会要辑稿》相关门类为本条来源起点；该条可固定编年位置，具体执行、规模和地方社会影响仍须保留来源差异。

\paragraph{1254年}\eventanchor{SS-1254-REGNAL-YEAR}{SS-1254-REGNAL-YEAR}　南宋以宋理宗在位第30年作为诏令、赋役与外交文书的纪年框架。账本所列《宋史》卷24—47〈本纪〉；《建炎以来系年要录》本年条；《宋会要辑稿》相关门类为本条来源起点；该条可固定编年位置，具体执行、规模和地方社会影响仍须保留来源差异。

\paragraph{1255年}\eventanchor{SS-1255-REGNAL-YEAR}{SS-1255-REGNAL-YEAR}　南宋以宋理宗在位第31年作为诏令、赋役与外交文书的纪年框架。账本所列《宋史》卷24—47〈本纪〉；《建炎以来系年要录》本年条；《宋会要辑稿》相关门类为本条来源起点；该条可固定编年位置，具体执行、规模和地方社会影响仍须保留来源差异。

\paragraph{1256年}\eventanchor{SS-1256-REGNAL-YEAR}{SS-1256-REGNAL-YEAR}　南宋以宋理宗在位第32年作为诏令、赋役与外交文书的纪年框架。账本所列《宋史》卷24—47〈本纪〉；《建炎以来系年要录》本年条；《宋会要辑稿》相关门类为本条来源起点；该条可固定编年位置，具体执行、规模和地方社会影响仍须保留来源差异。

\paragraph{1257年}\eventanchor{SS-1257-REGNAL-YEAR}{SS-1257-REGNAL-YEAR}　南宋以宋理宗在位第33年作为诏令、赋役与外交文书的纪年框架。账本所列《宋史》卷24—47〈本纪〉；《建炎以来系年要录》本年条；《宋会要辑稿》相关门类为本条来源起点；该条可固定编年位置，具体执行、规模和地方社会影响仍须保留来源差异。

\paragraph{1258年}\eventanchor{SS-1258-EVENT-141}{SS-1258-EVENT-141}　蒙哥发动大规模攻宋计划。账本所列《宋史》卷24—47〈本纪〉；《建炎以来系年要录》本年条；《宋会要辑稿》相关门类为本条来源起点；该条可固定编年位置，具体执行、规模和地方社会影响仍须保留来源差异。

\paragraph{1258年}\eventanchor{SS-1258-REGNAL-YEAR}{SS-1258-REGNAL-YEAR}　南宋以宋理宗在位第34年作为诏令、赋役与外交文书的纪年框架。账本所列《宋史》卷24—47〈本纪〉；《建炎以来系年要录》本年条；《宋会要辑稿》相关门类为本条来源起点；该条可固定编年位置，具体执行、规模和地方社会影响仍须保留来源差异。

\paragraph{1259年}\eventanchor{SS-1259-EVENT-142}{SS-1259-EVENT-142}　钓鱼城战事与蒙哥去世改变战争进程。账本所列《宋史》卷24—47〈本纪〉；《建炎以来系年要录》本年条；《宋会要辑稿》相关门类为本条来源起点；该条可固定编年位置，具体执行、规模和地方社会影响仍须保留来源差异。

\paragraph{1259年}\eventanchor{SS-1259-REGNAL-YEAR}{SS-1259-REGNAL-YEAR}　南宋以宋理宗在位第35年作为诏令、赋役与外交文书的纪年框架。账本所列《宋史》卷24—47〈本纪〉；《建炎以来系年要录》本年条；《宋会要辑稿》相关门类为本条来源起点；该条可固定编年位置，具体执行、规模和地方社会影响仍须保留来源差异。

\paragraph{1260年}\eventanchor{SS-1260-REGNAL-YEAR}{SS-1260-REGNAL-YEAR}　南宋以宋理宗在位第36年作为诏令、赋役与外交文书的纪年框架。账本所列《宋史》卷24—47〈本纪〉；《建炎以来系年要录》本年条；《宋会要辑稿》相关门类为本条来源起点；该条可固定编年位置，具体执行、规模和地方社会影响仍须保留来源差异。

\paragraph{1261年}\eventanchor{SS-1261-REGNAL-YEAR}{SS-1261-REGNAL-YEAR}　南宋以宋理宗在位第37年作为诏令、赋役与外交文书的纪年框架。账本所列《宋史》卷24—47〈本纪〉；《建炎以来系年要录》本年条；《宋会要辑稿》相关门类为本条来源起点；该条可固定编年位置，具体执行、规模和地方社会影响仍须保留来源差异。

\paragraph{1262年}\eventanchor{SS-1262-REGNAL-YEAR}{SS-1262-REGNAL-YEAR}　南宋以宋理宗在位第38年作为诏令、赋役与外交文书的纪年框架。账本所列《宋史》卷24—47〈本纪〉；《建炎以来系年要录》本年条；《宋会要辑稿》相关门类为本条来源起点；该条可固定编年位置，具体执行、规模和地方社会影响仍须保留来源差异。

\paragraph{1263年}\eventanchor{SS-1263-REGNAL-YEAR}{SS-1263-REGNAL-YEAR}　南宋以宋理宗在位第39年作为诏令、赋役与外交文书的纪年框架。账本所列《宋史》卷24—47〈本纪〉；《建炎以来系年要录》本年条；《宋会要辑稿》相关门类为本条来源起点；该条可固定编年位置，具体执行、规模和地方社会影响仍须保留来源差异。

\paragraph{1264年}\eventanchor{SS-1264-EVENT-143}{SS-1264-EVENT-143}　理宗去世，度宗即位，贾似道权势上升。账本所列《宋史》卷24—47〈本纪〉；《建炎以来系年要录》本年条；《宋会要辑稿》相关门类为本条来源起点；该条可固定编年位置，具体执行、规模和地方社会影响仍须保留来源差异。

\paragraph{1264年}\eventanchor{SS-1264-REGNAL-YEAR}{SS-1264-REGNAL-YEAR}　南宋以宋理宗在位第40年作为诏令、赋役与外交文书的纪年框架。账本所列《宋史》卷24—47〈本纪〉；《建炎以来系年要录》本年条；《宋会要辑稿》相关门类为本条来源起点；该条可固定编年位置，具体执行、规模和地方社会影响仍须保留来源差异。

\paragraph{1265年}\eventanchor{SS-1265-REGNAL-YEAR}{SS-1265-REGNAL-YEAR}　南宋以宋度宗在位第1年作为诏令、赋役与外交文书的纪年框架。账本所列《宋史》卷24—47〈本纪〉；《建炎以来系年要录》本年条；《宋会要辑稿》相关门类为本条来源起点；该条可固定编年位置，具体执行、规模和地方社会影响仍须保留来源差异。

\paragraph{1266年}\eventanchor{SS-1266-REGNAL-YEAR}{SS-1266-REGNAL-YEAR}　南宋以宋度宗在位第2年作为诏令、赋役与外交文书的纪年框架。账本所列《宋史》卷24—47〈本纪〉；《建炎以来系年要录》本年条；《宋会要辑稿》相关门类为本条来源起点；该条可固定编年位置，具体执行、规模和地方社会影响仍须保留来源差异。

\paragraph{1267年}\eventanchor{SS-1267-REGNAL-YEAR}{SS-1267-REGNAL-YEAR}　南宋以宋度宗在位第3年作为诏令、赋役与外交文书的纪年框架。账本所列《宋史》卷24—47〈本纪〉；《建炎以来系年要录》本年条；《宋会要辑稿》相关门类为本条来源起点；该条可固定编年位置，具体执行、规模和地方社会影响仍须保留来源差异。

\paragraph{1268年}\eventanchor{SS-1268-EVENT-144}{SS-1268-EVENT-144}　蒙古军开始围攻襄阳、樊城。账本所列《宋史》卷24—47〈本纪〉；《建炎以来系年要录》本年条；《宋会要辑稿》相关门类为本条来源起点；该条可固定编年位置，具体执行、规模和地方社会影响仍须保留来源差异。

\paragraph{1268年}\eventanchor{SS-1268-REGNAL-YEAR}{SS-1268-REGNAL-YEAR}　南宋以宋度宗在位第4年作为诏令、赋役与外交文书的纪年框架。账本所列《宋史》卷24—47〈本纪〉；《建炎以来系年要录》本年条；《宋会要辑稿》相关门类为本条来源起点；该条可固定编年位置，具体执行、规模和地方社会影响仍须保留来源差异。

\paragraph{1269年}\eventanchor{SS-1269-REGNAL-YEAR}{SS-1269-REGNAL-YEAR}　南宋以宋度宗在位第5年作为诏令、赋役与外交文书的纪年框架。账本所列《宋史》卷24—47〈本纪〉；《建炎以来系年要录》本年条；《宋会要辑稿》相关门类为本条来源起点；该条可固定编年位置，具体执行、规模和地方社会影响仍须保留来源差异。

\paragraph{1270年}\eventanchor{SS-1270-REGNAL-YEAR}{SS-1270-REGNAL-YEAR}　南宋以宋度宗在位第6年作为诏令、赋役与外交文书的纪年框架。账本所列《宋史》卷24—47〈本纪〉；《建炎以来系年要录》本年条；《宋会要辑稿》相关门类为本条来源起点；该条可固定编年位置，具体执行、规模和地方社会影响仍须保留来源差异。

\paragraph{1271年}\eventanchor{SS-1271-REGNAL-YEAR}{SS-1271-REGNAL-YEAR}　南宋以宋度宗在位第7年作为诏令、赋役与外交文书的纪年框架。账本所列《宋史》卷24—47〈本纪〉；《建炎以来系年要录》本年条；《宋会要辑稿》相关门类为本条来源起点；该条可固定编年位置，具体执行、规模和地方社会影响仍须保留来源差异。

\paragraph{1272年}\eventanchor{SS-1272-REGNAL-YEAR}{SS-1272-REGNAL-YEAR}　南宋以宋度宗在位第8年作为诏令、赋役与外交文书的纪年框架。账本所列《宋史》卷24—47〈本纪〉；《建炎以来系年要录》本年条；《宋会要辑稿》相关门类为本条来源起点；该条可固定编年位置，具体执行、规模和地方社会影响仍须保留来源差异。

\paragraph{1273年}\eventanchor{SS-1273-EVENT-145}{SS-1273-EVENT-145}　襄阳、樊城失守。账本所列《宋史》卷24—47〈本纪〉；《建炎以来系年要录》本年条；《宋会要辑稿》相关门类为本条来源起点；该条可固定编年位置，具体执行、规模和地方社会影响仍须保留来源差异。

\paragraph{1273年}\eventanchor{SS-1273-REGNAL-YEAR}{SS-1273-REGNAL-YEAR}　南宋以宋度宗在位第9年作为诏令、赋役与外交文书的纪年框架。账本所列《宋史》卷24—47〈本纪〉；《建炎以来系年要录》本年条；《宋会要辑稿》相关门类为本条来源起点；该条可固定编年位置，具体执行、规模和地方社会影响仍须保留来源差异。

\paragraph{1274年}\eventanchor{SS-1274-EVENT-146}{SS-1274-EVENT-146}　元军大举南下，南宋朝廷陷入危机。账本所列《宋史》卷24—47〈本纪〉；《建炎以来系年要录》本年条；《宋会要辑稿》相关门类为本条来源起点；该条可固定编年位置，具体执行、规模和地方社会影响仍须保留来源差异。

\paragraph{1274年}\eventanchor{SS-1274-REGNAL-YEAR}{SS-1274-REGNAL-YEAR}　南宋以宋度宗在位第10年作为诏令、赋役与外交文书的纪年框架。账本所列《宋史》卷24—47〈本纪〉；《建炎以来系年要录》本年条；《宋会要辑稿》相关门类为本条来源起点；该条可固定编年位置，具体执行、规模和地方社会影响仍须保留来源差异。

\paragraph{1275年}\eventanchor{SS-1275-EVENT-147}{SS-1275-EVENT-147}　南宋军在丁家洲等战场失利，贾似道被罢黜。账本所列《宋史》卷24—47〈本纪〉；《建炎以来系年要录》本年条；《宋会要辑稿》相关门类为本条来源起点；该条可固定编年位置，具体执行、规模和地方社会影响仍须保留来源差异。

